%!TEX root = ../main/main.tex

The definitions of documents, spans, mappings, regex formulas, and
variable-set automata below follow Section~2 of Florenzano et
al.~\cite{FlorenzanoRUVV20}.
We present the models in the order VA, eVA, and logical-VA, making our
conventions for runs and captures explicit. The logical-VA follows
Riveros et al.~\cite{RiverosJV23}. These preliminaries provide the
vocabulary needed to study a Glushkov construction with capture
variables.

\paragraph{Notation and conventions.}
We use $\gamma$ for a regex formula with capture variables (an RGX) and
$E$ for an ordinary regular expression over an alphabet $\Delta$. The
expression $\widehat{E}$ is obtained from $E$ by replacing every symbol
occurrence with its position; it is the \emph{marked expression} used in
the Glushkov construction. We write $G_\gamma$ for the Glushkov
variable-set automaton associated with $\gamma$
(Section~\ref{subsec:glushkov-captures}). The symbol $+$ denotes union,
$\cdot$ denotes concatenation, and $*$ denotes the Kleene star. Ordinary
expressions may contain $\varnothing$; regex formulas do not, following
the syntax of Florenzano et al. We write $\operatorname{Pos}(E)$ for the
positions of an expression and $\operatorname{pos}(d)$ for the positions
of a document. When counting states, we always restrict to states
reachable from the initial state. The Glushkov construction is first
formulated for VAs and then lifted to the eVA and logical-VA models.

\subsection{Documents and spans}\label{subsec:documents-spans}

Let $\Sigma$ be a finite alphabet. Let $\Sigma^*$ denote the set of finite
words over $\Sigma$, and let $\varepsilon$ denote the empty word.

\begin{definition}[Document and positions]\label{def:document}
A \emph{document} is a word $d=a_1\cdots a_n\in\Sigma^*$.
Its length is $|d|=n$, and its set of positions is
\[
  \operatorname{pos}(d)=\{1,\ldots,n+1\}.
\]
Position $i\leq n$ lies immediately before the character $a_i$;
position $n+1$ lies after the last character. In particular,
$\operatorname{pos}(\varepsilon)=\{1\}$.
\end{definition}

\begin{definition}[Span]\label{def:span}
A \emph{span} of $d$ is a half-open interval $s=[i,j\rangle$ with
$1\leq i\leq j\leq |d|+1$. The set of spans of $d$ is
\[
  \operatorname{span}(d)
  =\{[i,j\rangle\mid 1\leq i\leq j\leq |d|+1\}.
\]
The content of $s$ is the substring
$d(s)=d(i,j)=a_i\cdots a_{j-1}$ when $i<j$, and
$\varepsilon$ when $i=j$. Its length is $j-i$.
For adjacent spans $s_1=[i,j\rangle$ and $s_2=[j,k\rangle$, their
concatenation 
is $s_1\cdot s_2=[i,k\rangle$.
\end{definition}

Distinct spans may have the same content. For instance, in
$d=\mathtt{aba}$, the spans $[1,2\rangle$ and $[3,4\rangle$ both
contain $\mathtt{a}$ but refer to different positions.

\subsection{Variables, mappings, and document spanners}\label{subsec:mappings}

Let $\mathcal V$ be a fixed set of variables, disjoint from $\Sigma$.
The set $\mathcal V$ need not be finite, although each mapping and each
automaton considered here uses only finitely many variables.

\begin{definition}[Mapping]\label{def:mapping}
A \emph{mapping over $d$} is a function
\[
  \mu:\operatorname{dom}(\mu)\longrightarrow\operatorname{span}(d),
  \qquad \operatorname{dom}(\mu)\subseteq\mathcal V\text{ finite}.
\]
Thus, a mapping is a partial assignment of variables to spans.
The \emph{empty mapping}, denoted by $\varnothing$, has empty domain.
We write $[x\mapsto s]$ for the mapping with domain $\{x\}$ that assigns
$s$ to $x$. Let $\operatorname{Map}(d)$ be the set of all mappings over $d$.
\end{definition}

\begin{definition}[Compatibility]\label{def:compatible-mappings}
Mappings $\mu_1,\mu_2\in\operatorname{Map}(d)$ are \emph{compatible},
written $\mu_1\sim\mu_2$, if
\[
  \mu_1(x)=\mu_2(x)
  \quad\text{for every }x\in
  \operatorname{dom}(\mu_1)\cap\operatorname{dom}(\mu_2).
\]
Their union $\mu_1\cup\mu_2$ is then the mapping with domain
$\operatorname{dom}(\mu_1)\cup\operatorname{dom}(\mu_2)$ that agrees
with each mapping on its domain.
\end{definition}

\begin{definition}[Document spanner]\label{def:spanner}
A \emph{document spanner} is a function $P$ assigning to every document
$d\in\Sigma^*$ a set $P(d)\subseteq\operatorname{Map}(d)$.
Each element of $P(d)$ is an extraction result.
\end{definition}

The result $P(d)=\varnothing$ contains no outputs, whereas
$P(d)=\{\varnothing\}$ contains one output: the mapping that assigns
no variables. An empty span $[i,i\rangle$ is a value that may be assigned
to a variable and is distinct from the empty mapping.

\begin{definition}[Regex formula]\label{def:rgx}
A \emph{regex formula} (RGX) over $\Sigma$ and $\mathcal V$ is an
expression generated by
\[
  \gamma ::= \varepsilon \mid a \mid x\{\gamma\} \mid \gamma\cdot\gamma
             \mid \gamma + \gamma \mid \gamma^{*},
\]
where $a$ ranges over $\Sigma$ and $x$ ranges over $\mathcal V$.
Write $\operatorname{var}(\gamma)$ for the set of variables occurring
in $\gamma$, and write $\mathrm{RGX}$ for the class of all regex formulas.
\end{definition}

The grammar does not include an explicit empty-language constant;
ordinary expressions, introduced in Section~\ref{subsec:tau}, do.

An RGX denotes a document spanner. Its semantics is the two-layer
relation of Florenzano et al., reproduced below. The first line below
defines the spanner $\llbracket\gamma\rrbracket_d$; the six clauses
after it define the relation $[\gamma]_d$. For a document $d$, a pair
$(s,\mu)$ belongs to $[\gamma]_d$ when $\gamma$ matches the substring
$d(s)$ and the match produces the mapping $\mu$. The spanner keeps only
the mappings whose match covers the whole document:
\[
  \llbracket\gamma\rrbracket_d
  =\{\mu\mid ([1,|d|+1\rangle,\mu)\in[\gamma]_d\}.
\]
The six clauses of the table are exactly the six forms of the grammar.

\vspace{1\baselineskip}
\noindent
\begin{minipage}{\linewidth}
\normalfont
\begin{center}
Table 2.\quad Two-layer semantics $[\gamma]_d$ and induced spanner
$\llbracket\gamma\rrbracket_d$ of an RGX $\gamma$ over a document $d$
\end{center}
\vspace{-0.5\baselineskip}
\hrule
\smallskip
\noindent{$\displaystyle
\begin{aligned}
  \llbracket\gamma\rrbracket_d
    &= \{\mu\mid ([1,|d|+1\rangle,\mu)\in[\gamma]_d\} \\
  [\varepsilon]_d
    &= \{(s,\varnothing)\mid s\in\mathrm{span}(d)\text{ and }d(s)=\varepsilon\} \\
  [a]_d
    &= \{(s,\varnothing)\mid s\in\mathrm{span}(d)\text{ and }d(s)=a\} \\
  [x\{\gamma\}]_d
    &= \{(s,\mu)\mid\exists(s,\mu')\in[\gamma]_d:
        x\notin\mathrm{dom}(\mu')\text{ and }\mu=[x\to s]\cup\mu'\} \\
  [\gamma_1\cdot\gamma_2]_d
    &= \{(s,\mu)\mid\exists(s_1,\mu_1)\in[\gamma_1]_d,
        \exists(s_2,\mu_2)\in[\gamma_2]_d: \\
    &\qquad s=s_1\cdot s_2,\ \mathrm{dom}(\mu_1)\cap\mathrm{dom}(\mu_2)
        =\varnothing,\text{ and }\mu=\mu_1\cup\mu_2\} \\
  [\gamma_1 + \gamma_2]_d
    &= [\gamma_1]_d\cup[\gamma_2]_d \\
  [\gamma^*]_d
    &= [\varepsilon]_d\cup[\gamma]_d\cup[\gamma^2]_d
       \cup[\gamma^3]_d\cup\cdots
\end{aligned}
$}
\smallskip
\hrule
\smallskip
\noindent Here, $\gamma^2$ is shorthand for $\gamma\cdot\gamma$, similarly
$\gamma^3$ for $\gamma\cdot\gamma\cdot\gamma$, and so forth.
\end{minipage}
\par\smallskip
\noindent Source: Florenzano et al.~\cite[Table~2]{FlorenzanoRUVV20}.
\par\smallskip

\newpage

\subsection{Variable-set automata}\label{subsec:va}

For each $x\in\mathcal V$, introduce two distinct formal symbols:
$x\vdash$ for \emph{opening} $x$, and $\dashv x$ for \emph{closing} $x$.
These symbols are called \emph{variable markers}. Their set is
\[
  \Gamma_{\mathcal V}
  =\{x\vdash,\dashv x\mid x\in\mathcal V\},
\]
and is disjoint from $\Sigma$.

\begin{definition}[VA]\label{def:va}
A \emph{variable-set automaton} (VA) over $\Sigma$ and $\mathcal V$ is
a tuple $A=(Q,q_0,F,\delta)$, where $Q$ is a finite set of states,
$q_0\in Q$ is the initial state, $F\subseteq Q$ is the set of final states,
and $\delta$ is a finite transition relation
\[
  \delta\subseteq Q\times(\Sigma\cup\Gamma_{\mathcal V})\times Q.
\]
Transitions $(q,a,q')$ with $a\in\Sigma$ are \emph{letter transitions}.
Transitions $(q,x\vdash,q')$ and $(q,\dashv x,q')$ are
\emph{variable transitions}. Define
\[
  \operatorname{var}(A)
  =\{x\in\mathcal V\mid
       (q,x\vdash,q')\in\delta\text{ or }(q,\dashv x,q')\in\delta
       \text{ for some }q,q'\in Q\}.
\]
\end{definition}

This definition does not require determinism. Variable transitions consume
no characters and retain their markers as labels; they are not
$\varepsilon$-labelled transitions.

\begin{definition}[Configuration and step]\label{def:configuration}
A \emph{configuration} of $A$ over $d$ is a pair $(q,i)$ with
$q\in Q$ and $i\in\operatorname{pos}(d)$. A step
$(q,i)\xrightarrow{o}_d(q',i')$ is allowed if $(q,o,q')\in\delta$ and
one of the following conditions holds:
\begin{enumerate}
  \item $o=a\in\Sigma$, $i\leq |d|$, $a_i=a$, and $i'=i+1$;
  \item $o\in\Gamma_{\mathcal V}$ and $i'=i$.
\end{enumerate}
\end{definition}

\begin{definition}[Run and acceptance of a VA]\label{def:va-run}
A \emph{run} of $A$ over $d$ is a finite sequence
\[
  \rho=(q_0,i_0)\xrightarrow{o_1}_d(q_1,i_1)
       \xrightarrow{o_2}_d\cdots
       \xrightarrow{o_m}_d(q_m,i_m),
\]
where every step satisfies Definition~\ref{def:configuration},
$i_0=1$, and $i_m=|d|+1$. A run with $m=0$ is allowed when these
conditions hold. The run is \emph{accepting} if $q_m\in F$.
\end{definition}

The positions $i_0,\ldots,i_m$ are non-decreasing. Each letter step
advances one position, whereas each variable step keeps the current
position. The endpoint conditions require a VA run to traverse the
entire document.

\begin{definition}[Valid run]\label{def:valid-run}
A run $\rho$ is \emph{valid} if, for every $x\in\mathcal V$, exactly
one of the following cases holds:
\begin{enumerate}
  \item neither $x\vdash$ nor $\dashv x$ occurs in $o_1\cdots o_m$;
  \item there are unique indices $r,s$ with $1\leq r<s\leq m$ such that
        $o_r=x\vdash$ and $o_s=\dashv x$, with no other occurrences
        of either marker.
\end{enumerate}
\end{definition}

Each variable used in a valid run is therefore opened once and closed once,
in that order. Captures of distinct variables need not be properly nested:
the marker sequence $x\vdash\,y\vdash\,\dashv x\,\dashv y$ is allowed.

\begin{definition}[Mapping induced by a run]\label{def:run-mapping}
A valid run $\rho$ induces a mapping $\mu_\rho$ whose domain consists
of exactly the variables whose markers occur in $\rho$.
If $o_r=x\vdash$ and $o_s=\dashv x$, then
\[
  \mu_\rho(x)=[i_{r-1},i_{s-1}\rangle.
\]
The endpoints are the document positions at which the markers are executed.
Since markers consume no characters, $i_{r-1}=i_r$ and $i_{s-1}=i_s$.
\end{definition}

The monotonicity of positions gives $i_{r-1}\leq i_{s-1}$, so that
$\mu_\rho(x)$ is a span of $d$. A valid run may induce an empty span
when the opening and closing markers are executed at the same position.

\begin{definition}[Semantics of a VA]\label{def:va-semantics}
The document spanner defined by $A$ is denoted by $\llbracket A\rrbracket$
and satisfies
\[
  \llbracket A\rrbracket_d
  =\{\mu_\rho\mid\rho\text{ is a valid accepting run of }A\text{ over }d\}.
\]
\end{definition}

The semantics is a set: if several runs induce the same mapping,
that mapping occurs only once. An accepting run that is not valid
contributes no output.

\begin{definition}[Sequentiality and functionality]\label{def:sequential-functional}
A valid run of $A$ is \emph{functional} if it uses all variables in
$\operatorname{var}(A)$; equivalently,
$\operatorname{dom}(\mu_\rho)=\operatorname{var}(A)$.
A VA is \emph{sequential} (sVA) if all its accepting runs are valid,
for every document. It is \emph{functional} (fVA) if all its accepting
runs are functional, for every document.
\end{definition}

Every fVA is sequential. Functionality does not require a unique extraction
result: it requires every accepting run to assign all variables validly.

%%%%

\subsection{From regex formulas to ordinary expressions}
\label{subsec:tau}

\begin{definition}[Ordinary regular expression]
\label{def:ordinary-regex}
Let $\Delta$ be a finite alphabet. An \emph{ordinary regular expression}
over $\Delta$ is generated by
\[
  E ::= \varnothing \mid \varepsilon \mid b \mid E+E \mid E\cdot E \mid E^{*},
\]
where $b$ ranges over $\Delta$. Its language
$\mathcal L(E)\subseteq\Delta^{*}$ is defined by
\begin{align*}
  \mathcal L(\varnothing)&=\varnothing, &
  \mathcal L(\varepsilon)&=\{\varepsilon\}, &
  \mathcal L(b)&=\{b\},\\
  \mathcal L(E_1+E_2)&=\mathcal L(E_1)\cup\mathcal L(E_2), &
  \mathcal L(E_1\cdot E_2)&=\mathcal L(E_1)\cdot\mathcal L(E_2), &
  \mathcal L(E^{*})&=\mathcal L(E)^{*}.
\end{align*}
\end{definition}

The alphabet of an ordinary expression may properly contain $\Sigma$.
Captures are eliminated by treating markers as letters of that larger
alphabet.

\begin{definition}[Translation of a regex formula]
\label{def:tau}
Let $\Omega=\Sigma\cup\Gamma_{\mathcal V}$.
The translation $\tau$ sends each regex formula to an ordinary regular
expression over $\Omega$:
\begin{align*}
  \tau(\varepsilon)&=\varepsilon, &
  \tau(a)&=a, &
  \tau(x\{\gamma\})&=x\vdash\cdot\tau(\gamma)\cdot\dashv x,\\
  \tau(\gamma_1\cdot\gamma_2)&=\tau(\gamma_1)\cdot\tau(\gamma_2), &
  \tau(\gamma_1+\gamma_2)&=\tau(\gamma_1)+\tau(\gamma_2), &
  \tau(\gamma^{*})&=\tau(\gamma)^{*}.
\end{align*}
Each marker is one symbol of $\Omega$. Thus $\tau$ is the homomorphism
of expressions that replaces a capture by the concatenation of its
opening marker, the translation of its body, and its closing marker.
\end{definition}

\begin{definition}[Valid word and induced mapping]
\label{def:word-mapping}
For $w\in\Omega^{*}$, write $\pi_{\Sigma}(w)$ for the word obtained by
erasing every marker. A word $w=o_1\cdots o_m\in\Omega^{*}$ is
\emph{valid} when, for every variable $x$, either neither marker of $x$
occurs in $w$, or there are unique indices $h,k$ with
$1\leq h<k\leq m$, $o_h=x\vdash$, and $o_k=\dashv x$.
A valid word induces the mapping $\mu_w$ whose domain consists of the
variables whose markers occur in $w$. If those markers occur at
indices $h<k$, then
\[
  \mu_w(x)=
  \bigl[1+|\pi_{\Sigma}(o_1\cdots o_{h-1})|,\;
        1+|\pi_{\Sigma}(o_1\cdots o_{k-1})|\bigr\rangle.
\]
\end{definition}

The indices $h$ and $k$ count symbols of $\Omega$, including markers.
The span endpoints count only symbols of $\Sigma$, starting at document
position $1$. An opening and a closing with no letter of $\Sigma$
between them induce an empty span.

\begin{definition}[Word semantics]
\label{def:word-semantics}
Let $E$ be an ordinary regular expression over $\Omega$. On a document
$d$ it denotes the set of mappings
\[
  \llbracket E\rrbracket_d
  =\{\mu_w\mid w\in\mathcal L(E),\ 
     w\text{ is valid, and }\pi_{\Sigma}(w)=d\}.
\]
\end{definition}

The table fixes $\llbracket\gamma\rrbracket_d$ from the regex formula.
Definition~\ref{def:word-semantics} fixes $\llbracket\tau(\gamma)\rrbracket_d$
from the ordinary expression. Both use the same brackets; the argument
shows which clause applies. Their equality is not claimed in this
section.

%%%

\subsection{The Glushkov construction}\label{subsec:glushkov}

We present the construction from
Giammarresi, Ponty, and Wood~\cite[Section~3.2]{GiammarresiPW98}.
Here, $E$ is an ordinary regular expression over a finite alphabet
$\Delta$, built from $\varnothing$, $\varepsilon$, and letters using
union, concatenation, and Kleene star. Its language is denoted by
$\mathcal L(E)$. In the application of Section~\ref{sec:results},
$\Delta=\Omega$.
This construction yields an automaton recognising $\mathcal L(E)$;
capture variables will require a separate extension.

\paragraph{Positions and the marked expression}
Number the occurrences of letters in $E$ from left to right, starting at
$1$. If there are $n$ occurrences, write
$\operatorname{Pos}(E)=\{1,\ldots,n\}$ and let
$\ell_E:\operatorname{Pos}(E)\to\Delta$ map each position to the letter
occurring there. Replacing each letter occurrence by its position gives
the marked expression $\widehat{E}$ over $\operatorname{Pos}(E)$. The symbols
$\varnothing$ and $\varepsilon$ are left unchanged and receive no positions.
These positions identify occurrences in the expression, rather than
positions in an input document.

\paragraph{First, last, and follow sets}
Define
\[
\begin{aligned}
  \operatorname{first}(E)
    &=\{i\in\operatorname{Pos}(E)\mid
        \exists u\in\operatorname{Pos}(E)^*: iu\in\mathcal L(\widehat{E})\},\\
  \operatorname{last}(E)
    &=\{i\in\operatorname{Pos}(E)\mid
        \exists u\in\operatorname{Pos}(E)^*: ui\in\mathcal L(\widehat{E})\},\\
  \operatorname{follow}(i,E)
    &=\{j\in\operatorname{Pos}(E)\mid
        \exists u,v\in\operatorname{Pos}(E)^*: uijv\in\mathcal L(\widehat{E})\},
        \qquad i\in\operatorname{Pos}(E).
\end{aligned}
\]
Thus, $\operatorname{first}(E)$ and $\operatorname{last}(E)$ contain the
positions that can begin and end a word of $\mathcal L(\widehat{E})$, respectively.
The set $\operatorname{follow}(i,E)$ contains the positions that can occur
immediately after $i$ in such a word.

\begin{samepage}
\begin{definition}[Glushkov automaton]\label{def:glushkov-automaton}
The \emph{Glushkov automaton} of $E$ is the nondeterministic finite automaton
\[
  M_E=(Q_E,\Delta,\delta_E,0,F_E),
\]
where $Q_E=\{0\}\cup\operatorname{Pos}(E)$, $0$ is the initial state,
and the transition relation is
\[
\begin{aligned}
  \delta_E={}&\{(0,\ell_E(j),j)\mid j\in\operatorname{first}(E)\}\\
             &{}\cup\{(i,\ell_E(j),j)\mid
                 i\in\operatorname{Pos}(E),\ j\in\operatorname{follow}(i,E)\}.
\end{aligned}
\]
The set of final states is
\[
  F_E=
  \begin{cases}
    \operatorname{last}(E)\cup\{0\}, & \text{if }\varepsilon\in\mathcal L(E),\\
    \operatorname{last}(E), & \text{otherwise}.
  \end{cases}
\]
\end{definition}
\end{samepage}

The initial state is included in $F_E$ when $E$ accepts the empty word.
This case is made explicit in
Br\"uggemann-Klein and Wood~\cite[Definition~2.3]{BruggemannKleinW98};
the abbreviated description in~\cite[Section~3.2]{GiammarresiPW98}
lists only $\operatorname{last}(E)$. The condition
$\varepsilon\in\mathcal L(E)$ is equivalent to
$\varepsilon\in\mathcal L(\widehat{E})$, because marking preserves
nullability. We state it on $E$ to emphasize that nullability is a
property of the unmarked expression.

The construction satisfies $\mathcal L(M_E)=\mathcal L(E)$.
It has $n+1$ states and no $\varepsilon$-transitions. No transition enters
the initial state, and every transition entering a position $j$ is
labelled by $\ell_E(j)$. There is at most $n+n^2$ transitions, giving
the quadratic worst-case transition bound stated
in~\cite[Section~3.2]{GiammarresiPW98}.
